% The spoken body of the homily for the First Sunday of Advent, said on Sunday
% 29 November 2026. Continuous preaching, ready to read aloud after the
% Gospel: no heading, no direction to the preacher, no citation inside the
% speech. The loci, the relation to the reviewed interpretations, the
% audience, the word count and the pace estimate stand in the terminal note,
% which no one reads out. The blank lines are the paragraph breaks a speaker
% recovers his place by. The speech tells the Gospel's scene and its two
% responses to the signs; gives St Gregory the Great's reason for the
% prophecy, his aging world and his reading of the lifted head; sets the
% Gradual's waiting beside the Gospel's expectation; gives Gregory's road and
% his haste, and St Paul's hour with St John Chrysostom's neighbours before
% dawn and his garment; tells St Augustine's sweet burden of the world, his
% "presently" and his deliverance;
% reads the Collect's petition and the Secret's cleansing; makes one response
% from the Introit and Offertory antiphon; and closes on Gregory's summer and
% the Gospel's command. The honorific is tied to the name so that "St" never
% ends a line.
\raggedcolumns

Our Lord spoke these words in the Temple at Jerusalem, a few days before his
Passion. He had just foretold that the city would be compassed with an army
and trodden down by the Gentiles. Then he spoke of signs in the sun and the
moon and the stars, of nations in distress at the roaring of the sea, and of
men ``withering away for fear and expectation of what shall come upon the
whole world.'' A moment later he told his disciples, ``When these things begin
to come to pass, look up and lift up your heads, because your redemption is at
hand.''

St~Gregory the Great preached on this Gospel to the people of Rome, and he
begins with the reason the Lord foretold such things. The Lord wishes to find
us ready, Gregory says, and so he tells what evils will come upon a world grown
old, ``to restrain us from loving it.'' He makes known the blows that come
before its end, so that if we will not fear God in quiet, we may at least fear
his judgment when we are worn down by blows. Gregory compares the world to a
human body. In youth the chest is strong and the neck thick; in old age the back
bends, the neck withers, the breath comes short, and even when no illness is
present, health itself is often a kind of sickness. The world flourished in
its early years, and now it is weighed down by its age. ``Do not love, my
brothers,'' he says, ``what you see cannot stand for long.''

Gregory also explains the command. To lift up your heads, he says, is to
gladden your hearts, ``because while the world of which you are not friends is
ending, the redemption you have sought draws near.'' Scripture often puts the
head for the mind, and to lift up the head is to raise our minds to the joys
of our home in heaven. Those who love God are told to rejoice at the world's
end, because they will soon find the one they love, while the one they did not
love passes away. Gregory's reason is a sentence of St~James: ``Whosoever
therefore will be a friend of this world becometh an enemy of God.'' Grief at
the world's ruin, he says, belongs to those who have planted the roots of the
heart in love of it.

So the men who wither and the disciples who look up see the same signs, and
what divides them is what they love. The Gradual prayed, ``None of them that
wait on thee shall be confounded,'' and its words come back in the
Offertory. In the Latin, the Gradual's word for those who wait comes from the
same verb as the Gospel's word for the expectation of the men who wither. Both
are waiting, for different things: the men in the Gospel expect what will
come upon the world they love, and the psalmist waits for God. Then the
Alleluia asked him, ``Shew us, O Lord, thy mercy; and grant us thy salvation.''
It can stand as the prayer of people who have sought the redemption that
Gregory says is drawing near.

Gregory puts a question to his people. ``What is mortal life but a road? And
consider, my brothers, what it is to grow weary in the labour of the road, and
yet not to want the road to end.'' Since we know the joys of our home in
heaven, he says, we ought to hasten to them as fast as we can, and wish to
arrive by a shorter road.

We can hear the same haste in the Epistle, where St~Paul tells the Romans,
``It is now the hour for us to rise from sleep. For now our salvation is
nearer than when we believed. The night is passed, and the day is at hand.''
St~John Chrysostom, preaching on these words, explains why Paul gives that
reason. Those who have newly believed are most eager at the start, and as time
goes on their earnestness is likely to wither away to nothing. Paul asks for
the reverse: the nearer the prize, the more wide awake a man ought to be, as
runners near the end of the course rouse themselves the more.

Chrysostom also describes what we all do before dawn. When the night is
pressing on toward morning and the swallow is heard, we wake our neighbour,
though it is still dark; and when the night is really going, we hurry one
another and say, ``It is day now!'' and set about the work of the day,
dressing and leaving our dreams. Paul goes on to the dressing. Cast off the
works of darkness, he says, and put on the armour of light; and at the last,
``put ye on the Lord Jesus Christ.'' Chrysostom observes that when Paul
speaks of sin he speaks of works, when he speaks of virtue he speaks of arms,
and at the end he gives us the Lord himself for a garment, so that Christ is
about us on every side and dwells within us.

A man can hear all this, know that it is the hour, and still not rise.
St~Augustine describes that state in his \work{Confessions}. He says the
things of the world weighed on him as sleep does, and he found the weight
sweet. He was like a man who wants to wake and sinks back into a heavy sleep,
who knows that it is time to get up and still lies there, displeased and yet
enjoying it. He heard God calling, ``Awake, thou that sleepest,'' and had
nothing to answer but drowsy words: ``Presently, lo, presently''; ``Leave me a
little while.'' ``But `presently, presently,' had no present.'' His mind was
held by the force of habit, even against its will, and he cried out with
St~Paul, ``Who shall deliver me from the body of this death?'' His answer was,
``Thy grace only, through Jesus Christ our Lord.'' When at last he was set
free, it was in reading these verses of Romans, the words that end today's
Epistle among them: ``put ye on the Lord Jesus Christ.''

At the beginning of this Mass the Church asked for that deliverance. The
Collect is addressed to the Son. It asks him to stir up his power and come,
``that, defended by Thee, we may deserve rescue from approaching dangers
brought on by our sins, and being set free by Thee, obtain our salvation.''
The danger it names lies in us, in what our sins have brought on, and the
deliverer it asks for is the one Augustine named. Paul tells us to rise, and
the Church asks Christ to come and set us free so that we can. We need to be
cleansed of our sins as well, and at the offering the Secret will ask that the
holy mysteries, ``cleansing us by their mighty power,'' make us come more pure
before God, who is their author.

The Mass began with a psalm: ``To thee have I lifted up my soul. In thee, O my
God, I put my trust.'' In the Latin, the psalm's ``lifted up'' is the verb of
the Gospel's ``lift up your heads,'' and by Gregory's account of that command
the psalmist who lifts up his soul is already doing what it asks. He does it
while his enemies threaten to laugh at him, as the disciples lift their heads
while the world is shaken. The same words return at the Offertory. When they
do, lift up to Christ whatever still holds you while you answer him
``presently,'' and ask him, as the Collect asked, to come and set you free
from it.

In the parable that follows the command, the fig tree puts out its fruit and
summer is near. Gregory reads the fruit of the world as its ruin, for the world
grows in order to fall. But the summer is the kingdom of God, he says,
``because then the clouds of our mourning pass, and the days of life shine
with the brightness of the eternal sun.'' So when these things begin to come
to pass, look up and lift up your heads, because your redemption is at hand.
